\section{Introduction}
\label{sec:intro}

Inductively coupled plasma (ICP) reactors are central to the
fabrication of semiconductor and microelectromechanical devices,
where they enable the controlled removal of material by reactive
neutrals generated in low-pressure electronegative discharges. The
design and operation of such reactors increasingly rely on
predictive models, both for early-stage process development and for
run-time control. Fully kinetic simulations---particle-in-cell,
hybrid Monte~Carlo, or two-dimensional fluid---remain too
computationally expensive to embed in optimization loops, digital
twins, or supervisory controllers, and the practical workhorse is
therefore a reduced-order model in which the chamber is described
by a zero- or one-dimensional set of balance equations supplemented
by phenomenological transport coefficients calibrated against a
small number of experimental observables
\citep{lieberman2005,chabert2011}.

The reduced-order modeling paradigm depends on a parameter
estimation step that is, in practice, an identification problem in
disguise. Fitted constants---wall recombination probabilities,
effective sticking coefficients, etch yields, residence times,
gas-heating coefficients---are adjusted to reproduce one or two
measured etch rates, and the resulting values are subsequently
interpreted as physical quantities and extrapolated to other
operating conditions. The interpretation implicitly assumes that
each fitted parameter has been pinned down by the data on its own
merits. In reality the parameters are coupled through the model
equations, and the data may constrain only certain combinations of
them, leaving directions in parameter space along which the fit is
degenerate. When the model also contains systematic errors---for
example, in the electron-impact rate coefficients used in the source
term---the degeneracies become a route by which those errors are
absorbed into the fitted values, distorting their physical
interpretation without producing any visible misfit at the
calibration points.

Sulfur hexafluoride (\sfsix{}) plasmas illustrate this difficulty in
a particularly sharp form. Their dominant low-energy electron loss
channel is dissociative attachment, which depletes the high-energy
tail of the electron energy distribution function (EEDF) below a
Maxwellian distribution at the same mean energy. Rate coefficients
for dissociative excitation and ionization, which are integrals of
the cross section against the EEDF tail, are therefore strongly
sensitive to the EEDF model adopted. Reduced ICP simulations
routinely substitute a Maxwellian distribution at an effective
electron temperature for the true EEDF when computing rate
coefficients from cross-section libraries. For \sfsix{} this
substitution is known to be poor, but what has not been
characterized quantitatively is how the resulting bias propagates
through a full calibration to the fitted depletion or transport
parameters of a reduced model, and in particular whether the bias is
observable through the calibration data themselves.

This work addresses that gap. We report a reduced-order \sfsix{} ICP
model in which the absorbed power, the electron density, the
dissociation source, and the radial fluorine transport are coupled
through a one-parameter feedstock-depletion correction with a
characteristic timescale $\beta$. The model is calibrated against
two literature etch-rate anchors at 200\,W and 2000\,W
\citep{mansano1999,panduranga2019}. With Maxwellian rate
coefficients computed from the Biagi \sfsix{} cross-section set
\citep{biagi}, the calibration returns $\beta \approx 14$\,ms,
almost two orders of magnitude shorter than any plausible gas
residence time for the target reactor class. We show that this is
not a failure of the depletion mechanism but a structural feature
of the calibration: through an analytical identifiability argument
we prove in the isothermal limit, and extend numerically under gas
heating, that two-anchor calibration constrains only the product
$\beta \cdot \kdiss \cdot \nel(P_{\mathrm{high}})$, and we verify
the result numerically by demonstrating that artificial rescaling
of $\kdiss$ shifts the fitted $\beta$ by exactly the inverse factor
while leaving every observable prediction unchanged. We then
exploit the identifiability theorem to make a falsifiable
prediction: replacing the Maxwellian rate coefficients with two-term
Boltzmann coefficients from BOLSIG$^+$ \citep{hagelaar2005} should
rescale $\beta$ by the ratio of the rate-source-dependent invariant.
The predicted shift,
$\beta_{\mathrm{Bolt}}/\beta_{\mathrm{Max}} \approx 35.7$, agrees
with the actual refit to better than $0.5\,\%$ once the secondary
shift in $\nel$ from the modified power balance is included. The
Boltzmann-calibrated $\beta = 509$\,ms lies within a nominal
geometric residence-time band for typical 10\,mTorr ICP chambers.

We further show that the same identifiability theorem implies a
strong observability constraint on EEDF inference from etch-rate
data. At the calibration $E/N$, the Maxwellian- and
Boltzmann-calibrated models yield etch-rate predictions
indistinguishable to the tolerance of the numerical fit across the
entire power sweep, while off the calibration $E/N$ the predictions
diverge by up to 85\,\%. A single etch-rate measurement at one
non-anchor $(E/N, x_{\mathrm{Ar}})$ configuration would distinguish
the two EEDF assumptions and lift the calibration degeneracy, with
implications for the design of multi-condition calibration
campaigns.

The remainder of the paper is organized as follows.
Section~\ref{sec:model} summarizes the reduced model and the
calibration framework. Section~\ref{sec:baseline} reports the
Maxwellian baseline calibration and the resulting $\beta$.
Section~\ref{sec:identifiability} contains the central
identifiability analysis and its numerical verification.
Section~\ref{sec:boltzmann} reports the Boltzmann-corrected
calibration and the comparison with the identifiability prediction.
Section~\ref{sec:observability} develops the observability result
and the corresponding experimental-design recommendations.
Section~\ref{sec:discussion} discusses the limits of the methodology
and the conclusions.
